\documentclass{article}
\usepackage[utf8]{inputenc}
\usepackage[T1]{fontenc}
\usepackage{hyperref}
\usepackage{url}
\usepackage{booktabs}
\usepackage{amsfonts}
\usepackage{nicefrac}
\usepackage{microtype}
\usepackage{amsmath,amssymb}
\usepackage{geometry}
\geometry{letterpaper, margin=1in}

\title{Supplementary Material: StateShift: Tracking State-Dependent Reasoning Recovery Across Post-Training}
\author{Anonymous Authors \\ NeurIPS 2027 Evaluations \& Datasets Track}

\begin{document}
\maketitle

\appendix

\section{Mathematical Formulations \& Estimand Proofs}

\subsection{Factorial Fixed-Effects Equivalence}
Consider the two-way fixed-effects linear probability model for rollout continuation outcome $y_{i,s,t,k}$:
\begin{equation}
y_{i,s,t,k} = \alpha_i + \beta_1 \mathbb{I}(s = R) + \beta_2 \mathbb{I}(t = 256) + \Gamma_{256} [\mathbb{I}(s = R) \times \mathbb{I}(t = 256)] + \varepsilon_{i,s,t,k}
\end{equation}
where $\alpha_i$ represents problem-level fixed effects, $s \in \{R, C\}$ is the experimental state condition, and $t \in \{0, 256\}$ is the checkpoint index.
Taking conditional expectations:
\begin{align}
\mathbb{E}[y_{i,R,256,k}] &= \alpha_i + \beta_1 + \beta_2 + \Gamma_{256} \\
\mathbb{E}[y_{i,R,0,k}]   &= \alpha_i + \beta_1 \\
\mathbb{E}[y_{i,C,256,k}] &= \alpha_i + \beta_2 \\
\mathbb{E}[y_{i,C,0,k}]   &= \alpha_i
\end{align}
Subtracting control deltas from recovery deltas:
\begin{equation}
(\mathbb{E}[y_{i,R,256,k}] - \mathbb{E}[y_{i,R,0,k}]) - (\mathbb{E}[y_{i,C,256,k}] - \mathbb{E}[y_{i,C,0,k}]) = (\beta_2 + \Gamma_{256}) - \beta_2 = \Gamma_{256}
\end{equation}
Because the experimental design evaluates every problem instance $i$ under all four condition-by-checkpoint cells with balanced sample size $K=16$, the OLS estimator for $\Gamma_{256}$ is algebraically identical to the sample double-difference of cell means:
\begin{equation}
\hat{\Gamma}_{256} = (\hat{\mu}_{R,256} - \hat{\mu}_{R,0}) - (\hat{\mu}_{C,256} - \hat{\mu}_{C,0}) = 0.320485 - 0.202919 = \mathbf{+0.117566}
\end{equation}

\section{Measurement Disattenuation and Rollout Sampling Noise}
\label{app:disattenuation}

In Section~\ref{sec:orthogonality}, we observed an empirical problem-level correlation between recovery gain $\Delta_{R,i}$ and clean gain $\Delta_{C,i}$ of $r = 0.0794$ ($p = 0.0910$), with empirical $R^2 = 0.63\%$. Here, we provide the formal derivation showing that finite rollout sampling ($K=16$) severely attenuates this observed correlation.

\subsection{Error-in-Variables and Sampling Variance Derivation}
Let $p_{i,s,t} \in [0, 1]$ denote the true, latent continuation success probability for problem instance $i \in \{1, \dots, N\}$, state condition $s \in \{R, C\}$, and checkpoint $t \in \{0, 256\}$.
Under stochastic decoding with $K$ i.i.d. rollouts per cell, the observed number of successes is $S_{i,s,t} \sim \text{Binomial}(K, p_{i,s,t})$, and the cell sample mean is $\hat{p}_{i,s,t} = S_{i,s,t} / K$.
The empirical problem-level change between baseline $t=0$ and post-training $t=256$ is:
\begin{equation}
\hat{\Delta}_{s,i} = \hat{p}_{i,s,256} - \hat{p}_{i,s,0} = \Delta_{s,i}^* + e_{s,i}
\end{equation}
where $\Delta_{s,i}^* = p_{i,s,256} - p_{i,s,0}$ represents the true latent change, and $e_{s,i} = (\hat{p}_{i,s,256} - p_{i,s,256}) - (\hat{p}_{i,s,0} - p_{i,s,0})$ represents estimation error due to finite rollout sampling.

Because rollouts at $t=0$ and $t=256$ are sampled independently, the sampling variance of $e_{s,i}$ conditional on the true probabilities is:
\begin{equation}
\sigma^2_{\text{sampling},s,i} = \text{Var}(e_{s,i} \mid p_{i,s}) = \frac{p_{i,s,256}(1 - p_{i,s,256})}{K} + \frac{p_{i,s,0}(1 - p_{i,s,0})}{K}
\end{equation}
By the law of total variance:
\begin{equation}
\text{Var}(\hat{\Delta}_{s,i}) = \text{Var}(\Delta_{s,i}^*) + \mathbb{E}\left[\sigma^2_{\text{sampling},s,i}\right]
\end{equation}
The classical reliability ratio $\text{Rel}(s) \in [0, 1]$ of the observed gain is defined as the proportion of observed variance attributable to true latent variation:
\begin{equation}
\text{Rel}(s) = \frac{\text{Var}(\Delta_{s,i}^*)}{\text{Var}(\hat{\Delta}_{s,i})} = 1 - \frac{\bar{\sigma}^2_{\text{sampling},s}}{\text{Var}(\hat{\Delta}_{s,i})}
\end{equation}

\subsection{Empirical Estimates from the Controlled Cohort}
Evaluating these terms across the $N=454$ problem blocks with $K=16$:
\begin{enumerate}
    \item \textbf{Recovery Condition ($s = R$)}:
    \begin{itemize}
        \item Total observed variance: $\text{Var}(\hat{\Delta}_{R,i}) = 0.02958$.
        \item Average binomial sampling variance: $\bar{\sigma}^2_{\text{sampling},R} = 0.02511$.
        \item Sampling noise proportion: $\frac{0.02511}{0.02958} = \mathbf{84.89\%}$.
        \item True latent variance: $\text{Var}(\Delta_{R,i}^*) = 0.02958 - 0.02511 = 0.00447$.
        \item Reliability: $\text{Rel}(R) = \mathbf{0.1511}$.
    \end{itemize}
    \item \textbf{Control Condition ($s = C$)}:
    \begin{itemize}
        \item Total observed variance: $\text{Var}(\hat{\Delta}_{C,i}) = 0.03396$.
        \item Average binomial sampling variance: $\bar{\sigma}^2_{\text{sampling},C} = 0.02727$.
        \item Sampling noise proportion: $\frac{0.02727}{0.03396} = \mathbf{80.30\%}$.
        \item True latent variance: $\text{Var}(\Delta_{C,i}^*) = 0.03396 - 0.02727 = 0.00669$.
        \item Reliability: $\text{Rel}(C) = \mathbf{0.1970}$.
    \end{itemize}
\end{enumerate}

\subsection{Spearman's Disattenuation Correction}
Because the rollout errors $e_{R,i}$ and $e_{C,i}$ arise from independent stochastic decodings, $\text{Cov}(e_{R,i}, e_{C,i}) = 0$. The observed sample covariance directly reflects the latent covariance: $\text{Cov}(\hat{\Delta}_{R,i}, \hat{\Delta}_{C,i}) = \text{Cov}(\Delta_{R,i}^*, \Delta_{C,i}^*)$.
Applying Spearman's classical disattenuation correction:
\begin{equation}
r_{\text{adj}} = \frac{r_{\text{obs}}}{\sqrt{\text{Rel}(R) \times \text{Rel}(C)}} = \frac{0.0794}{\sqrt{0.1511 \times 0.1970}} = \frac{0.0794}{0.1725} \approx \mathbf{0.4605}
\end{equation}
The disattenuated shared variance is $R^2_{\text{adj}} \approx (0.4605)^2 = \mathbf{21.21\%}$.
This establishes that while problem-level recovery improvements are uncoupled in raw empirical measurements ($R^2 = 0.63\%$), finite rollout sampling accounts for over $80\%$ of observed variance, and the true latent constructs share moderate overlap ($R^2 \approx 21\%$).

\section{Checkpoint Trajectory and Order-Restricted Isotonic Modeling}
\label{app:trajectory}

To investigate the temporal progression of state-conditioned recovery, we tracked $\Gamma_t$ across nine checkpoints spanning the DeepScaleR post-training trajectory.

\subsection{Checkpoint Sample Sizes and Observed Outcomes}
Table~\ref{tab:app_trajectory} presents the complete numerical records across all nine evaluated checkpoints. At the confirmatory endpoints ($t=0$ and $t=256$), evaluations were conducted with $K=16$ rollouts per problem cell ($N=454, 7,264$ rollouts per condition). Intermediate checkpoints ($t \in \{32, 64, \dots, 224\}$) were evaluated with $K=2/3$ rollouts per cell ($908$ to $1,362$ rollouts per condition).

\begin{table}[h]
\centering
\caption{Trajectory of Recovery and Control Success Rates, Unconstrained Interaction Contrast ($\hat{\Gamma}_t$), and Isotonic PAVA-Fitted Values ($\tilde{\Gamma}_t$) Across Checkpoints.}
\label{tab:app_trajectory}
\small
\begin{tabular}{rcccccc}
\toprule
Checkpoint $t$ & Rollouts / Cell & $\mu_{R,t}$ & $\mu_{C,t}$ & Observed $\hat{\Gamma}_t$ & $95\%$ Blocked CI & Isotonic PAVA $\tilde{\Gamma}_t$ \\
\midrule
0   & $K=16$ & $0.3834$ & $0.3892$ & $+0.0000$ & $[+0.0000, +0.0000]$ & $+0.0000$ \\
32  & $K=2/3$ & $0.4681$ & $0.4406$ & $+0.0333$ & $[+0.0023, +0.0677]$ & $+0.0333$ \\
64  & $K=2/3$ & $0.4890$ & $0.4611$ & $+0.0337$ & $[+0.0041, +0.0645]$ & $+0.0337$ \\
96  & $K=2/3$ & $0.5705$ & $0.4989$ & $+0.0774$ & $[+0.0461, +0.1084]$ & $+0.0673$ \\
128 & $K=2/3$ & $0.5969$ & $0.5279$ & $+0.0748$ & $[+0.0435, +0.1068]$ & $+0.0673$ \\
160 & $K=2/3$ & $0.6068$ & $0.5528$ & $+0.0598$ & $[+0.0278, +0.0911]$ & $+0.0673$ \\
192 & $K=2/3$ & $0.6553$ & $0.5635$ & $+0.0976$ & $[+0.0674, +0.1292]$ & $+0.0963$ \\
224 & $K=2/3$ & $0.6696$ & $0.5804$ & $+0.0950$ & $[+0.0631, +0.1265]$ & $+0.0963$ \\
256 & $K=16$ & $0.7039$ & $0.5921$ & $+0.1176$ & $[+0.0958, +0.1403]$ & $+0.1176$ \\
\bottomrule
\end{tabular}
\end{table}

\subsection{Pool Adjacent Violators Algorithm (PAVA)}
To model the overall monotonic progression without imposing a rigid linear or log-linear functional form, we apply the Pool Adjacent Violators Algorithm (PAVA). Let $\hat{\Gamma}_t$ denote the unconstrained estimate at checkpoint step $t_j$ for $j \in \{1, \dots, 9\}$. PAVA solves the weighted order-restricted least-squares problem:
\begin{equation}
\min_{\tilde{\Gamma}_1 \le \tilde{\Gamma}_2 \le \dots \le \tilde{\Gamma}_9} \sum_{j=1}^9 w_j (\hat{\Gamma}_{t_j} - \tilde{\Gamma}_{t_j})^2
\end{equation}
Because $\hat{\Gamma}_{128} = 0.0748$ and $\hat{\Gamma}_{160} = 0.0598$ violate monotonicity relative to $\hat{\Gamma}_{96} = 0.0774$, PAVA pools these three steps into a single plateau with weighted average value $\tilde{\Gamma} = 0.0673$. Similarly, the dip at step 224 ($\hat{\Gamma}_{224} = 0.0950$ vs. $\hat{\Gamma}_{192} = 0.0976$) is pooled into a plateau with value $\tilde{\Gamma} = 0.0963$.

The total sum of squared errors between the empirical estimates and the isotonic fit is:
\begin{equation}
\text{SSE} = \sum_{j=1}^9 (\hat{\Gamma}_{t_j} - \tilde{\Gamma}_{t_j})^2 = 0.000602
\end{equation}
The close fit ($\text{SSE} = 0.0006$) demonstrates that the macro-trajectory is well summarized by a non-decreasing profile. However, we clarify that PAVA is an exploratory smoothing technique that constructs a monotonic sequence by algorithm design; it does not constitute an inferential proof of step-by-step monotonicity in the underlying training process.

\section{Decontamination Protocol \& Operator Catalog}
Solutions from MATH-500 were segmented into distinct reasoning deductions. To ensure structural recoverability and invalidity:
\begin{enumerate}
    \item \textbf{OP\_CONSTANT\_PERTURB ($N=161$)}: Modifies an intermediate arithmetic constant while preserving equation structure (e.g., $3x + 4 = 19 \implies 3x = 12$).
    \item \textbf{OP\_SIGN\_FLIP ($N=143$)}: Inverts a sign in an algebraic manipulation (e.g., $x^2 - 5x + 6 = 0 \implies x^2 + 5x + 6 = 0$).
    \item \textbf{OP\_FRACTION\_FLIP ($N=84$)}: Inverts a fraction within a denominator or summation (e.g., $\frac{1}{n^3} \implies \frac{n^3}{1}$).
\end{enumerate}

\section{Verification Oracles}
All evaluations use a deterministic Python verification harness coupling SymPy 1.13 with a LaTeX boxed expression extractor. The oracle verifies symbolic equivalence by expanding and simplifying algebraic differences:
\begin{equation}
\text{Verify}(A_{\text{pred}}, A^*) = \mathbb{I}\left(\text{simplify}(A_{\text{pred}} - A^*) = 0\right)
\end{equation}
Zero neural evaluators or prompt-based judges were utilized.

\end{document}
